% ============================================================
% base/templates/format.tex   内置格式配置（基线 · 仓库里的源）
%
% 这一层只放"长什么样"：页面、颜色、字体、字号、页标题、间距。
% 不含任何组件（组件在 components.tex）。
%
% 【这份文件不参与 deck 编译】
%   新建 deck 时只复制 `style.tex`、`main.tex`、`pages/`。
%   本文件与 components.tex 留在 skill 里，是给 deck**抄**的源：
%   deck 的 style.tex 自带等价内容，并另加载本 deck 的 .slide-forge/ 覆盖。
%   —— 不复制它们进 deck，是为了避免与
%   <deck>/slide/.slide-forge/templates/format.tex 同名共存，
%   那种"两个同名文件都得改"的坑迟早会把内容改分叉。
%
% 【三层架构里各自的位置】
%   base/templates/style.tex          内置基线 → 复制进 deck，唯一的事实来源
%   base/templates/format.tex|components.tex   本文件，仓库里的基线源
%   <deck>/slide/.slide-forge/templates/       本 deck 的本地层（自动生成）
%   evolved/templates/format.tex|components.tex 吸收来的仓库（不参与编译）
%
% ★ 不要直接改本文件来调格式。要改请写进对应 deck 的
%   .slide-forge/templates/format.tex，满意后再走「吸收」进 evolved/。
% ============================================================

% ============================================================
% 一、页面与导航
% ============================================================
\usetheme{default}
\usefonttheme{professionalfonts}
\usepackage{tabularx}   % 表格用它让表宽自适应
\usepackage{graphicx}   % 图片页
\setbeamertemplate{navigation symbols}{}
\setbeamertemplate{headline}{}
\setbeamertemplate{footline}{}
\setbeamersize{text margin left=0.90cm,text margin right=0.90cm}

% ============================================================
% 二、颜色
% ------------------------------------------------------------
% 正文一律黑色，页标题用强调色；全篇不使用斜体
% （中文没有真斜体字形，合成倾斜会发虚）。
% 要强调某个词时用 \textbf 加粗，一句话里不要超过一处。
% ============================================================
\definecolor{SlideInk}{HTML}{14161A}       % 正文：近黑（唯一正文色）
\definecolor{SlideAccent}{HTML}{26308F}    % 页标题：低饱和深靛蓝
\definecolor{SlideBG}{HTML}{FBFAF7}        % 暖白背景

\setbeamercolor{background canvas}{bg=SlideBG}
\setbeamercolor{normal text}{fg=SlideInk,bg=SlideBG}
\setbeamercolor{frametitle}{fg=SlideAccent,bg=SlideBG}
\setbeamercolor{itemize item}{fg=SlideInk}
\setbeamercolor{itemize subitem}{fg=SlideInk}
\setbeamercolor{structure}{fg=SlideAccent}

% ============================================================
% 三、字体
% ------------------------------------------------------------
% 中文：微软雅黑　西文与数字：Segoe UI
% 微软雅黑没有真斜体，所以全篇不使用 \emph / \textit。
% 可用字重只有 Regular 与 Bold 两档 —— 正好对应"正文 / 强调"。
% ============================================================
\setCJKsansfont{Microsoft YaHei}
\setsansfont{Segoe UI}
\renewcommand{\familydefault}{\sfdefault}

% ============================================================
% 四、字号与行距：全篇只有两级
% ------------------------------------------------------------
% 页标题档 / 正文档。正文里的一切（含姓名、单位、要点）都用
% \SlideTextFont，不存在第三级字号。
% 每一级的第二个数字就是行距，不再另外设 \linespread。
% ============================================================
\newcommand{\SlideTitleFont}{\fontsize{20}{26}\selectfont}  % 页标题档
\newcommand{\SlideTextFont}{\fontsize{12}{17}\selectfont}   % 正文档（唯一）

% 封面档：只用于封面与致谢页的主标题。
% 全篇唯一的例外 —— 封面上的姓名如果和正文一样大就不成其为封面，
% 但它只出现在"封面 / 致谢"这两页，正文页一律不用。
\newcommand{\SlideCoverFont}{\fontsize{26}{32}\selectfont}

\setbeamerfont{frametitle}{size=\SlideTitleFont,series=\bfseries}
\setbeamerfont{normal text}{size=\SlideTextFont,series=\mdseries}
\setbeamerfont{itemize/enumerate body}{size=\SlideTextFont,series=\mdseries}

% ============================================================
% 五、页标题
% ------------------------------------------------------------
% 左上角，无分隔线、无页脚、无装饰。整份 deck 靠它对齐。
% ============================================================
\setbeamertemplate{frametitle}{%
  \nointerlineskip%
  \vspace*{0.50cm}%
  \begin{beamercolorbox}[wd=\paperwidth,leftskip=0.90cm,rightskip=0.90cm,ht=0.72cm,dp=0.06cm]{frametitle}%
    \usebeamerfont{frametitle}\insertframetitle%
  \end{beamercolorbox}%
}

% 页标题与正文之间的统一间距：整页呼吸感的起点。
% 必须用 \vspace*（带星号）：普通 \vspace 出现在段落起点时会被
% TeX 在换页处丢弃，结果就是"设了间距却没生效"。
\newlength{\SlideBodyTopSkip}
\setlength{\SlideBodyTopSkip}{0.40cm}

% ============================================================
% 六、行距与垂直节奏
% ============================================================
\setlength{\parindent}{0pt}
\setlength{\parskip}{0pt}

% 列表不用符号，改用悬挂缩进：某条折行时第二行与首行文字左对齐，
% 条目边界依然清楚。圆点在 12pt 下要么发灰、要么与基线打架。
\setlength{\leftmargini}{0pt}
\setbeamertemplate{itemize item}{}
\setbeamertemplate{itemize subitem}{}
\setbeamertemplate{itemize/enumerate body begin}{\SlideTextFont}
